%! Linear Regression — Unit 1, Lesson 1.5 — Lesson Plan
% Gradual-release plan (course-wide shape, 2026-09-01): warm-up → guided notes (with guided
% practice) → individual practice → debrief → close. There is no group activity; the "you do" is
% done alone. No tiers, no exit ticket; homework is in-packet.
\documentclass[10pt]{article}
\usepackage{algebra2-article}
\usepackage{algebra2-boxes}
\usepackage{graphicx}   % -article does NOT load graphicx; needed for thumbnails/screenshots
\graphicspath{{images/}}
\newcommand{\TallMath}[1]{$\displaystyle #1\rule[-1.4em]{0pt}{3.2em}$}
\newcommand{\CourseName}{Algebra 2}
\newcommand{\MeetingLength}{60 minutes}
\newcommand{\UnitNumberName}{Unit 1: Linear Functions \quad}
\newcommand{\LessonNumberName}{Lesson 1.5: Linear Regression}

\begin{document}
\begin{center}
    {\Large\bfseries \CourseName \\
    {\small\bfseries \UnitNumberName \LessonNumberName}}
\end{center}

\begin{tcolorbox}[colback=forestbg, colframe=forest, boxrule=0.9pt, arc=2mm,
    left=2mm, right=2mm, top=1.5mm, bottom=1.5mm]
    \textbf{Primary Objective:} Students move from \emph{exact} rules to \emph{real, scattered
    data}. Given a \textbf{scatterplot} of bivariate data, they describe the \textbf{association}
    by \emph{direction}, \emph{form}, and \emph{strength}, and interpret the \textbf{correlation
    coefficient} $r$ --- its \emph{sign} is the direction and its \emph{size} the strength of the
    linear fit. They read a \textbf{line of best fit} $\hat y=ax+b$ produced with technology,
    \textbf{interpret its slope and $y$-intercept in context} with units, and \textbf{predict}
    with it, distinguishing \emph{interpolation} (inside the data) from \emph{extrapolation}
    (outside it) and judging whether a prediction is reasonable. They explain why
    \emph{correlation does not prove causation}. All scatterplots and lines are pre-drawn and all
    regression values are \emph{given} --- students read, interpret, predict, and justify.
    \\[2pt]
    \textbf{Standards (2023 VA SOL):} \textbf{A2.ST.2} --- apply the data cycle with a focus on
    \emph{bivariate data}: represent bivariate data with a scatterplot using technology
    (\textbf{A2.ST.2c}); decide whether a linear model is appropriate (\textbf{A2.ST.2d});
    determine the equation of the line of best fit using technology (\textbf{A2.ST.2e}); use the
    \textbf{correlation coefficient} to designate goodness of fit (\textbf{A2.ST.2f}); make
    predictions and critical judgments from the model (\textbf{A2.ST.2g}); and evaluate the
    reasonableness of the model in context (\textbf{A2.ST.2h}). This lesson closes Unit~1's
    linear strand.
    \\[2pt]
    \textbf{Lesson model:} \emph{Gradual release} --- I do, we do, you do alone. A warm-up
    activates prior knowledge; \textbf{Guided Notes} build the vocabulary with the class on one
    worked context (hours practised against free throws made) and close with a \emph{Guided
    Practice} box (the ``we do''); an \textbf{Individual Practice} block on the same page is the
    ``you do'' --- three problems, worked silently and alone; a \textbf{debrief} consolidates and
    surfaces the errors; \textbf{homework} extends the practice. \emph{There is no group activity
    in this lesson} --- today's independent work is done individually.
\end{tcolorbox}

\renewcommand{\arraystretch}{1.4}
\begin{skillbox}[Priority Ideas \& Skills]{goldbox}
\begin{tabularx}{\linewidth}{@{}X|X@{}}
\textbf{Priority Ideas \& Skills} & \textbf{Key Understandings (the ``why'')} \\[2pt]
\hline
\vspace{-6pt}
\begin{itemize}[leftmargin=1.1em, nosep, after=\vspace{2pt}]
  \item Describe a scatterplot by \textbf{direction}, \textbf{form}, and \textbf{strength}.
  \item Read $r$: sign $=$ direction, size $=$ strength; match $r$ to a plot.
  \item Interpret the slope of $\hat y=ax+b$ as a \emph{per-unit rate} and the intercept as a
        \emph{baseline}, both with units.
  \item Predict by substitution; label it interpolation or extrapolation.
  \item Judge a prediction against the story; reject an impossible one.
  \item Refuse a causal claim from correlation alone; name a lurking variable.
\end{itemize}
&
\vspace{-6pt}
\begin{itemize}[leftmargin=1.1em, nosep, after=\vspace{2pt}]
  \item Real data rarely lands on one line, so we summarise the \emph{trend} rather than find an
        exact rule --- the first time in Unit~1 the data is messy.
  \item \textbf{Target misconception:} that a negative $r$ is a \emph{weak} one. The sign is
        direction only; $r=-0.9$ is a stronger fit than $r=0.4$.
  \item The line runs through the \emph{middle} of the points, not through them; technology
        picks the line whose predictions are closest to the data overall.
  \item The line does not know the story's limits: outside the data it keeps going ($20$ made out
        of $15$ shots) while reality does not.
\end{itemize}
\end{tabularx}
\end{skillbox}

\begin{skillbox}[{Vocabulary, Concepts, \& Theorems}]{sky}
\renewcommand{\arraystretch}{1.3}
\begin{tabularx}{\linewidth}{@{}l X@{}}
\textbf{Scatterplot / bivariate data} & Paired measurements $(x,y)$ on the same individuals, plotted as points. \\
\textbf{Association} & The pattern in a scatterplot: \emph{direction} (positive / negative / none), \emph{form} (linear?), \emph{strength} (how tightly?). \\
\textbf{Correlation coefficient $r$} & A number in $[-1,1]$: its \emph{sign} is the direction, its \emph{size} $|r|$ the strength of the \emph{linear} relationship. \\
\textbf{Line of best fit} & The regression line $\hat y=ax+b$ from technology; $\hat y$ is a \emph{predicted} value. Slope $=$ per-unit rate; intercept $=$ baseline. \\
\textbf{Interpolation / extrapolation} & Predicting \emph{inside} the data's $x$-range vs.\ \emph{outside} it, where the trend may break. \\
\textbf{Lurking variable} & A hidden third variable that drives both measured ones --- why correlation does not prove causation. \\
\end{tabularx}
\end{skillbox}

\begin{fixedskillbox}[Lesson at a Glance --- \MeetingLength]{forestbg}
\renewcommand{\arraystretch}{1.25}
\begin{tabularx}{\linewidth}{@{}l c X X@{}}
\textbf{Phase} & \textbf{Min} & \textbf{Students} & \textbf{Teacher} \\
\hline
Warm-Up & 5 & Three prior skills, alone then check with a neighbour & Circulate; debrief item~1 aloud (``no line goes through all four points --- but do they mostly rise?'') \\
Guided Notes & 20 & Fill the vocabulary and the four sections with the class, then \emph{Guided Practice} together & \emph{I do / we do}: run the whole block on \emph{one} context, the free-throw data; two \texttt{work} blocks at the board plus the Guided Practice box \\
Individual Practice & 15 & \emph{Alone and quiet}: three problems --- match $r$ to plots, interpret a model, find where a line breaks & Circulate; questions, cues, prompts --- \emph{not} answers. Note who to call on in the debrief \\
Debrief & 10 & Correct their own pages in a second colour & Work all three individual-practice problems on the board; land ``sign is direction, size is strength'' and ``the line does not know the story'' \\
Close \& Homework & 10 & Start the homework page in class; preview the Unit~1 test and Unit~2 & Launch item~1 aloud, then circulate; name what changed today \\
\end{tabularx}
\end{fixedskillbox}

\begin{fixedskillbox}[Warm-Up --- Activate Prior Knowledge \& Spiral Review]{forestbg}
The Warm-Up rehearses exactly what today leans on. \textbf{(1)} Reading four labelled points off
a grid is the skill a scatterplot demands --- every scatterplot is just many points like these ---
and the follow-up (\emph{do they lie on one line? do they mostly rise?}) is the seed of
``trend, not rule.'' \textbf{(2)} Evaluating $y=4x+10$ at three inputs \emph{is} prediction: the
same arithmetic as substituting into $\hat y=ax+b$ in section~4. \textbf{(3)} Slope and intercept
from $y=ax+b$, the direction a positive or negative slope gives, and ``each time $x$ goes up by
$1$, $y$ goes up by\dots'' --- the exact sentence section~3 turns into a rate with units.
\textbf{Debrief item 1 aloud, briefly}, and land only this: \emph{no single line goes through all
four points, but they mostly rise --- hold on to the difference between a rule and a trend.}
\textbf{This page is set at 12pt} --- larger type than the rest of the packet, deliberately.
\end{fixedskillbox}

\begin{skillbox}[Hook]{forestbg}
    ``A coach records, for each player, the hours practised this week and the free throws made out
    of 15.'' Put the scatterplot on the board and ask: \textbf{``As hours go up, what do the free
    throws do?''} (They tend to go up.) \textbf{``Could a single line pass through every
    point?''} (No.) \textbf{``Could one line summarise the trend?''} (Yes.) Then the detail that
    carries the lesson: \textbf{``The player at 5 hours made 11; the player at 6 hours made 11.
    Does more practice \emph{guarantee} more baskets?''} (No --- it \emph{tends} to.) Say:
    \emph{no exact rule fits messy data, but a line of best fit can summarise the trend --- and
    today we read it, measure how well it fits, and predict with it.}
\end{skillbox}

\begin{skillbox}[Guided Notes --- the Lesson (20 min)]{forestbg}
\begin{multicols}{2}
\textbf{1. Scatterplots \& association.} Three pre-drawn plots: (a) positive and linear, (b)
negative, (c) no trend. Describe each three ways --- direction, form, strength --- with \emph{no
numbers yet}. Then classify the free-throw data the same way (positive, linear, strong).

\textbf{2. The correlation coefficient $r$.} One number in $[-1,1]$: sign is direction, size is
strength; the labelled number line. Match $0.90$, $-0.85$, $0.05$ to plots (a), (b), (c). Then the
two-column table, $r=0.9$ beside $r=-0.9$ --- same strength, opposite direction --- and the gold
caution: \emph{which is the stronger fit, $r=-0.9$ or $r=0.4$?} $|{-0.9}|=0.9>0.4$. \emph{The sign
is direction only.} This is the target misconception, and this box is where it dies.

\columnbreak
\textbf{3. The line of best fit.} Technology gives $\hat y=1.5x+2$, $r=0.97$, drawn through the
middle of the points. \textbf{Interpret in context}: the slope is a rate --- \emph{each additional
hour is associated with about $1.5$ more free throws} --- and the intercept is the baseline at
$0$ hours. Insist on units every time.

\textbf{4. Predicting --- and its limits.} Two \texttt{work} blocks side by side: $x=6$ gives
$11$ (inside $1$--$8$: interpolation, reasonable); $x=12$ gives $20$ (outside: extrapolation) ---
and the gold caution: only $15$ shots were attempted, so the line has left the story. Close with
ice cream and drownings: a lurking variable (hot weather) pushes both, so correlation is not
causation.

\textbf{Guided Practice (the ``we do'').} Exercise minutes against resting heart rate,
$\hat y=-0.4x+78$, $r=-0.88$: negative \emph{and} strong; the slope with units; $\hat y(30)=66$ in
a \texttt{work} block; interpolation; and ``does a strong $r$ prove causation?'' Do this one
\emph{with} the class --- it is the last thing they see before working alone, and it is the first
\emph{negative} association they interpret.
\end{multicols}
\end{skillbox}

\begin{skillbox}[Individual Practice --- \emph{On Your Own} (15 min, silent)]{redbox}
\textbf{Launch (1 min):} ``Same moves you just watched --- now nobody helps you. Three problems.
Direction from the sign, strength from the size; every slope gets its units; check every
prediction against the data's range. Stuck? Look up the page at the notes, not at your
neighbour.'' The block sits on the last page of the Guided Notes, so nothing new is handed out.
\begin{multicols}{2}
\textbf{What students do.} \textbf{Problem~1} matches $0.95$, $0.10$, $-0.70$ to three plots ---
$P$ the tight riser, $Q$ the shapeless cloud, $R$ the loose faller --- and names the strongest
fit ($P$). \textbf{Problem~2} interprets the salary model $\hat y=2.5x+38$ (salary in thousands):
the slope as \$2{,}500 \emph{per year}, the intercept as \$38{,}000, $\hat y(10)=63$ in a
\texttt{work} block read as \$63{,}000 and labelled interpolation, and $x=30$ labelled
extrapolation. \textbf{Problem~3} is the crux: hot cocoa against temperature,
$\hat y=-3x+260$ over $20$--$60^\circ$; the $40^\circ$ prediction is $140$ cups, but the
$90^\circ$ one is $-10$ --- impossible --- and they must say why in two lines. Together the three
cover \emph{read $r$}, \emph{interpret and predict}, and \emph{judge}.

\columnbreak
\textbf{What the teacher does.} Circulate the whole 15 minutes; keep it silent. Give a
\emph{question, cue, or prompt}, never an answer:
\begin{itemize}[leftmargin=1.1em, nosep]
  \item 1: ``Which plot falls? So which $r$ must be negative? Now how tightly?''
  \item 2: ``$2.5$ what, per what? The $y$-axis is in thousands.''
  \item 2: ``Is $10$ years inside the workers you measured? Is $30$?''
  \item 3: ``Can you sell $-10$ cups? So what has the line done at $90^\circ$?''
  \item 3: ``Where does the line hit zero cups? Is that inside the data?''
  \item Finished early: ``In problem~3, the caf\'e adds days from $60$ to $80^\circ$ and $r$
        drops to $-0.75$. What might the new points look like?''
\end{itemize}
Note names as you go: one student for each of the three problems, to put work on the board in the
debrief.
\end{multicols}
\end{skillbox}

\boxguard
\begin{skillbox}[Debrief (10 min --- what goes on the board, in a second color)]{forestbg}
Students correct their own Individual Practice in a second colour as you go. Work all three, but
protect the time for items~2 and~4.
\begin{enumerate}[leftmargin=1.3em, nosep, itemsep=2pt]
  \item \textbf{Problem~1 --- read $r$ in two steps.} Direction first: only $R$ falls, so only
        $R$ can take $-0.70$. Then strength: $P$ hugs its line ($0.95$), $Q$ has no line to hug
        ($0.10$). Ask who put $-0.70$ on $Q$ ``because it looks messy'' --- messy is about size,
        not sign.
  \item \textbf{The must-land moment: sign is direction, size is strength.} Write $r=-0.9$ and
        $r=0.4$ side by side and ask which points sit closer to their line. Point back at the gold
        box in notes section~2. Say it out loud: \emph{the minus sign tells you which way the
        trend runs; the $0.9$ tells you how tightly.}
  \item \textbf{Problem~2 --- units, or it is not an interpretation.} ``$2.5$'' is wrong;
        ``\$2{,}500 more salary per additional year of experience'' is right. $\hat y(10)=63$ is
        \$63{,}000, interpolation; $x=30$ is twice the longest career measured --- extrapolation.
  \item \textbf{Problem~3 --- the line does not know the story.} $\hat y(90)=-10$ cups. The line
        crosses zero at about $87^\circ$, far outside $20$--$60^\circ$. What to conclude instead:
        \emph{very hot days sell almost none} --- the model's job ends where the data ends.
\end{enumerate}
If time is short, cut item~1's discussion of $Q$ and protect items~2 and~4.
\end{skillbox}

\boxguard
\begin{skillbox}[Active Monitoring --- Watch For]{redbox}
\begin{itemize}[leftmargin=1.2em, nosep]
  \item \textbf{calling a negative $r$ ``weak''} (notes~2, Individual Practice~1, homework~2
        and~6) --- the lesson's target misconception. Send them to $|r|$, not to the sign;
  \item interpreting a slope with no units or no context --- ``the slope is $2.5$'' (Guided
        Practice~2, Individual Practice~2, homework~3);
  \item missing that salary is in \emph{thousands}, so $\hat y=63$ means \$63{,}000
        (Individual Practice~2);
  \item trusting an extrapolation because the arithmetic was right --- $20$ made out of $15$,
        $-10$ cups, $85$ cm on day $100$ (notes~4, Individual Practice~3, homework~3c);
  \item calling a tight \emph{curved} pattern ``$r\approx1$'' --- $r$ measures a linear fit only
        (homework~4b);
  \item claiming causation from a strong $r$, or being unable to name a lurking variable
        (Guided Practice~4, homework~5);
  \item order-of-operations slips substituting into $\hat y=ax+b$, especially with a negative
        slope (Guided Practice~4, Individual Practice~3).
\end{itemize}
Cold-call prompts: ``Direction, form, strength --- which is which here?'' ``Is this $r$ closer to
$-1$, $0$, or $1$, and how do you know?'' ``What does this slope mean in the story, with units?''
``Inside the data or outside --- and is that prediction reasonable?'' ``Does this prove one thing
\emph{causes} the other?''
\end{skillbox}

\boxguard
\begin{skillbox}[Reinforcement \& Extension]{goldbox}
\textbf{Homework} (in the packet, collected next class): matching $0.98$, $-0.60$, $0.05$ to three
plots and naming each direction; a deliberate \emph{sign-versus-size} pair ($r=-0.85$ against
$r=0.45$) with the ``why is negative not weak'' question; a model read from a \emph{data table}
(seedling height, $\hat y=0.8x+5$) --- interpret with units, an interpolation checked against the
table, and an extrapolation to day $100$ to reject; the special values $r=1$, $r=0$, and an
impossible $r=1.3$; the shoe-size / reading-level causation classic with a ``hold age fixed''
concept check; and an SOL-style multiple-choice item on the strongest of four $r$-values.
\textbf{Extension:} two candidate lines for five points, compared by total error size --- an
informal look at what technology minimises. \textbf{DeltaMath override:} interpreting $r$ and
predicting from a given line are well covered there --- swap in a DeltaMath set if you prefer
auto-graded practice; the reasonableness and causation items are thinner there, so keep
homework~3(c) and~5 if you do. \textbf{Preview:} this closes Unit~1. Next is the \textbf{Unit~1
Test}; then \textbf{Unit~2, Quadratic Functions}, where the data can curve and a \emph{parabola},
not a line, fits best.
\end{skillbox}

% ── Teacher notes — one per component, in packet order ──
\begin{teachernote}[Warm-Up]
Item~1 is the reading skill a scatterplot demands, and its two follow-ups are the whole lesson in
miniature: the four points do \emph{not} lie on a line, yet they \emph{mostly} rise --- a trend
without a rule. Item~2 \emph{is} prediction: substituting into $y=4x+10$ is exactly what
section~4 does with $\hat y=1.5x+2$; keep order of operations tight, because the Guided Practice
has a negative slope. Item~3 seeds the day's language --- a positive slope means $y$ rises with
$x$ (a positive association), and ``each time $x$ goes up by $1$, $y$ goes up by $3$'' is the
sentence section~3 turns into a rate with units. Debrief item~1 in one sentence and leave it.
\textbf{This page is set at 12pt} --- larger type than the rest of the packet, deliberately. It is
still one page; if you add to it, check that it stays one page.
\end{teachernote}

\begin{teachernote}[Guided Notes]
Twenty minutes, paced by section: 3 for the vocabulary and hook, 3 for section~1, 5 for
section~2, 4 for section~3, 3 for section~4, and 2 for the Guided Practice box worked \emph{with}
the class. \textbf{Keep the whole notes block on one context} --- the free-throw data --- so the
same eight points carry description, $r$, the line, and the prediction. The organising arc is
\emph{describe $\to$ measure $\to$ model $\to$ predict $\to$ judge}. Section~2 is the one to
protect: the $r=0.9$ / $r=-0.9$ table and its gold caution ($|{-0.9}|=0.9>0.4$) are where
``negative means weak'' dies. Section~4's two side-by-side \texttt{work} blocks make the
interpolation/extrapolation contrast visible on one line of the page; the ``$20$ out of $15$''
caution is the memorable moment. Finding $r$ and the equation is a technology task
(A2.ST.2c, e, f); every value here is \emph{given} so class time goes to reading and reasoning.
The stated $\hat y=1.5x+2$ and $r=0.97$ are what technology actually returns for the eight points
drawn --- if a student checks on a calculator, the numbers hold. If you fall behind, compress
section~1's plot (c) to one sentence --- but do \emph{not} cut section~4's extrapolation block;
Individual Practice~3 depends on it.
\end{teachernote}

\begin{teachernote}[Individual Practice]
Fifteen minutes, and \textbf{keep it silent} --- this is the only block all period where a student
finds out what they can do without help. Budget roughly 4, 5, and 6 minutes. The three problems
are deliberately the three \emph{jobs} of the lesson: read $r$ (problem~1), interpret and predict
(problem~2), and judge (problem~3). A student who matches problem~1 correctly but writes ``$2.5$
more'' with no units in problem~2 has the number and not the meaning --- send them to the
$y$-axis label, not to a rule. Watch for $-0.70$ landing on the shapeless plot $Q$ ``because it
is messy'' (messy is size, not sign), for $\hat y=63$ left as \$63 instead of \$63{,}000, and for
the $-10$ cups written down without comment --- the arithmetic is right; the judgment is the
point. Fast finishers reason about what new hot-day points would do to $r$; that is the work you
want on the board in the debrief.
\end{teachernote}

\begin{teachernote}[Homework]
Item~2 is the target misconception in the wild --- the two-line explanation should say ``the sign
is direction, the size is strength, and $0.85>0.45$.'' Item~3 is the interpretation core
(A2.ST.2g, h) read off a \emph{table} rather than a plot: part (a) demands units (cm per day),
part (b) is a safe interpolation the student can check against the table's $20$ cm, and part (c)
is the extrapolation that breaks --- $85$ cm on day $100$, when growth slows. Item~4(c) is the
boundary: $r=1.3$ is impossible, not ``very strong.'' Item~5(b) is the concept check --- holding
age fixed removes the lurking variable, so the correlation should collapse; a student who says
``still strong'' has not understood what a lurking variable does. Item~6 is the formative check:
sort into ``chose B (got it) / chose C (read negative as weak --- the target misconception, still
alive) / chose A or D (guessed)'' to decide whether the Unit~1 review opens with a one-minute
re-look at ``sign versus size.'' Two \texttt{work} blocks (items 3b and 3c). With the group
activity cut, the last 10 minutes are for \emph{starting} this page in class --- launch item~1
aloud, then circulate; the rest is due next class.
\end{teachernote}

\end{document}
